\begin{definitionbox}{Definition (Left Identity)}
\begin{definition}[Left Identity]
\label{def:left-identity}
\[
\forall x,\ e_L\mathbin{O}x=x.
\]
\end{definition}
\end{definitionbox}

\begin{remark*}[Standard quantified statement]
\[
e_L\in A
\quad\text{and}\quad
\forall x\in A,\ e_L\mathbin{O}x=x.
\]
\end{remark*}

\begin{remark*}[Predicate reading]
\[
\operatorname{LeftIdentity}(e_L,A,O).
\]
\end{remark*}

\begin{remark*}[Interpretation]
A left identity is tied to both the carrier and the operation; the same element
need not be a left identity for another operation.
\end{remark*}

\begin{dependencies}
\begin{itemize}
  \item \hyperref[def:binary-arithmetic-operation]{Binary Arithmetic Operation}
\end{itemize}
\end{dependencies}

\begin{definitionbox}{Definition (Right Identity)}
\begin{definition}[Right Identity]
\label{def:right-identity}
\[
\forall x,\ x\mathbin{O}e_R=x.
\]
\end{definition}
\end{definitionbox}

\begin{remark*}[Standard quantified statement]
\[
e_R\in A
\quad\text{and}\quad
\forall x\in A,\ x\mathbin{O}e_R=x.
\]
\end{remark*}

\begin{remark*}[Predicate reading]
\[
\operatorname{RightIdentity}(e_R,A,O).
\]
\end{remark*}

\begin{remark*}[Interpretation]
A right identity is an operation-relative property on the chosen carrier.
\end{remark*}

\begin{dependencies}
\begin{itemize}
  \item \hyperref[def:binary-arithmetic-operation]{Binary Arithmetic Operation}
\end{itemize}
\end{dependencies}

\begin{definitionbox}{Definition (Two-Sided Identity)}
\begin{definition}[Two-Sided Identity]
\label{def:two-sided-identity}
An element $e$ is a \emph{two-sided identity} for $O$ if it is both a left
identity and a right identity.
\end{definition}
\end{definitionbox}

\begin{remark*}[Standard quantified statement]
\[
\operatorname{LeftIdentity}(e,A,O)
\quad\text{and}\quad
\operatorname{RightIdentity}(e,A,O).
\]
\end{remark*}

\begin{remark*}[Predicate reading]
\[
\operatorname{Identity}(e,A,O).
\]
\end{remark*}

\begin{remark*}[Interpretation]
Two-sided identity packages the left and right identity laws for one operation
on one carrier.
\end{remark*}

\begin{dependencies}
\begin{itemize}
  \item \hyperref[def:left-identity]{Left Identity}
  \item \hyperref[def:right-identity]{Right Identity}
\end{itemize}
\end{dependencies}

\begin{definitionbox}{Definition (Absorbing Element)}
\begin{definition}[Absorbing Element]
\label{def:absorbing-element}
An element $z$ is an \emph{absorbing element} for a binary operation $O$ if
\[
z\mathbin{O}x=z
\quad\text{and}\quad
x\mathbin{O}z=z
\]
for every element $x$ in the domain of the operation.
\end{definition}
\end{definitionbox}

\begin{remark*}[Standard quantified statement]
\[
z\in A
\quad\text{and}\quad
\forall x\in A,\ z\mathbin{O}x=z \land x\mathbin{O}z=z.
\]
\end{remark*}

\begin{remark*}[Predicate reading]
\[
\operatorname{AbsorbingElement}(z,A,O).
\]
\end{remark*}

\begin{remark*}[Interpretation]
An absorbing element fixes the operation output whenever it appears in either
input slot.
\end{remark*}

\begin{dependencies}
\begin{itemize}
  \item \hyperref[def:binary-arithmetic-operation]{Binary Arithmetic Operation}
\end{itemize}
\end{dependencies}

\begin{lemmabox}{Lemma (Left and Right Identities Coincide)}
\begin{lemma}[Left and Right Identities Coincide]
\label{lem:left-right-identities-coincide}
If $e_L$ is a left identity and $e_R$ is a right identity for the same
operation, then $e_L=e_R$.
\hyperref[prf:left-right-identities-coincide]{\textit{Go to proof.}}
\end{lemma}
\end{lemmabox}

\begin{remark*}[Standard quantified statement]
\[
\operatorname{LeftIdentity}(e_L,A,O)
\land
\operatorname{RightIdentity}(e_R,A,O)
\Longrightarrow e_L=e_R.
\]
\end{remark*}

\begin{remark*}[Predicate reading]
\[
\operatorname{LeftIdentity}(e_L,A,O)
\land
\operatorname{RightIdentity}(e_R,A,O)
\Longrightarrow e_L=e_R.
\]
\end{remark*}

\begin{remark*}[Interpretation]
Left and right identities cannot disagree once they belong to the same carrier
and operation.
\end{remark*}

\begin{dependencies}
\begin{itemize}
  \item \hyperref[def:left-identity]{Left Identity}
  \item \hyperref[def:right-identity]{Right Identity}
\end{itemize}
\end{dependencies}

\begin{theorembox}{Theorem (Uniqueness of Identity)}
\begin{theorem}[Uniqueness of Identity]
\label{thm:uniqueness-of-identity}
A two-sided identity for a binary operation is unique.
\hyperref[prf:uniqueness-of-identity]{\textit{Go to proof.}}
\end{theorem}
\end{theorembox}

\begin{remark*}[Standard quantified statement]
\[
\operatorname{Identity}(e,A,O)\land\operatorname{Identity}(e',A,O)
\Longrightarrow e=e'.
\]
\end{remark*}

\begin{remark*}[Predicate reading]
\[
\operatorname{Unique}(\lambda u.\operatorname{Identity}(u,A,O)).
\]
\end{remark*}

\begin{remark*}[Interpretation]
Identity is unique as a property of a fixed carrier-operation pair.
\end{remark*}

\begin{dependencies}
\begin{itemize}
  \item \hyperref[lem:left-right-identities-coincide]{Left and Right Identities Coincide}
\end{itemize}
\end{dependencies}
